\newpage
\thispagestyle{empty}
\renewcommand{\abstractname}{\kaishu \sihao \textbf{摘要}}
\begin{abstract}
\noindent
这次编译原理上机实验由成员A和成员B分工完成。小题1中，成员A以带有普通宏、带参宏和条件编译的 C++ 阶乘程序为例，使用 GCC 和 Clang 观察预处理、词法分析、tree、CFG、RTL、汇编、目标文件和链接后的可执行文件；成员B以包含宏、数组和循环的 C17 程序为例，使用 Clang/LLVM 保存预处理结果、token、AST、CFG、LLVM IR、AArch64 汇编以及 ELF 符号和重定位信息，并比较 \texttt{-O0} 与 \texttt{-O2} 的输出差异。两部分的命令、截图和中间文件均保留在仓库中。

小题2中，两人共同设计了一个 SysY 程序：读入整型数组，计算偶数元素之和，并统计正数元素个数。成员A编写了等价的 LLVM IR，成员B编写了等价的 RISC-V 汇编，两种实现都链接 SysY 运行时库完成验证。LLVM IR 分别在本机程序和 \texttt{lli} 中测试，RISC-V 程序在 QEMU 中测试；三组统一数据的结果均与预期一致，RISC-V 版本补充的 100 组特殊和普通数据也全部通过。实验结果表明，LLVM IR 可以用基本块和 SSA 值表示循环、分支和数组访问，而 RISC-V 汇编还需要直接处理寄存器、栈帧和数组地址。

进阶部分以 $2\times4$、$4\times3$ 矩阵乘法观察 MLIR 中端处理。成员A整理了 \texttt{linalg.matmul} 到三层 \texttt{scf.for} 的表示变化，成员B核对基本块形式、LLVM dialect、LLVM IR 和 x86-64 目标文件。保存的 \texttt{matmul.ll} 能重新编译为导出 \texttt{matmul} 符号的 ELF 可重定位文件。

\noindent\textbf{关键词：} GCC；编译流程；SysY；LLVM IR；RISC-V；运行时库
\end{abstract}